\documentclass[a4paper,12pt]{article}

\usepackage{amsmath}
\usepackage{amssymb}
\usepackage{geometry}
\usepackage{graphicx}
\usepackage{siunitx}
\usepackage{hyperref}
\usepackage{caption}
\usepackage{float}

\geometry{margin=1in}
\sisetup{per-mode=symbol}

\title{Cookbook TS741}
\author{}
\date{}

\begin{document}

\maketitle
\vspace{-1em}
\tableofcontents
\newpage

\section{Overview}

The cookbook is a simple file that contains all the essential information about the pedal. If you want to recreate it while modifying the sound or changing some parameters, you can refer to this file and easily adjust resistors or capacitors to achieve the desired result.

\section{Input Buffer}

\begin{figure}[H]
    \centering
    \includegraphics[width=0.6\textwidth]{tube-screamer-input-buffer.png}
    \caption{Input buffer}
    \label{fig:input-buffer}
\end{figure}

The input buffer preserves the signal as much as possible by providing high input impedance and a unity gain.

\subsection{Input Impedance}

The input impedance can be calculated using the hybrid-$\pi$ model for the transistor:
\begin{equation}
    Z_{\text{in}} = R_1 + \left( R_2 \parallel \left[ r_\pi + (\beta + 1) R_3 \right] \right)
\end{equation}

Since the transistor's $\beta$ parameter is high, the input impedance approaches the value of resistor $R_2$. For proper operation, it is recommended to interface the circuit with a guitar having an impedance of at least \SI{500}{\kilo\ohm}.

\subsubsection{Voltage Gain}

The input buffer gain is approximately unity ($A_v \approx 1$), as its primary function is to preserve the signal without significant amplification.

\section{Clipping Stage}

\begin{figure}[H]
    \centering
    \includegraphics[width=0.6\textwidth]{tube-screamer-clipping-amplifier.png}
    \caption{Clipping stage}
    \label{fig:clipping-stage}
\end{figure}

The clipping stage is responsible for the characteristic distortion of the Tube Screamer. The gain of the operational amplifier in a non-inverting configuration is given by:
\begin{equation}
    A_v = 1 + \frac{R_6 + P_1}{R_4}
\end{equation}

where $P_1$ is the gain control potentiometer. Since high-value logarithmic potentiometers are not easily available, an effective method to increase the gain is to modify resistor $R_6$.

\section{Output Buffer}

\begin{figure}[H]
    \centering
    \includegraphics[width=0.6\textwidth]{tube-screamer-output-buffer.png}
    \caption{Output buffer}
    \label{fig:output-buffer}
\end{figure}

The output buffer reduces the circuit's output impedance, making it suitable for driving subsequent stages without signal degradation.

The output impedance is given by:
\begin{equation}
    Z_{\text{out}} = R_C \parallel \left( R_B + \left( R_{13} \parallel \frac{r_\pi + R_{12}}{\beta + 1} \right) \right)
\end{equation}

To ensure a low output impedance, it is advisable to choose a low value for resistor $R_C$.

\end{document}
